% class_diagram.tex  — TikZ source for the ComplexGitSync class overview figure.
%
% This file is \input{} from docs/architecture.tex.  It draws four related
% diagrams on separate pages so each fits comfortably within an A4 margin:
%
%   Figure 1 — Class overview by tier (Tier 1/2/3 headline classes, plus the
%              Document-layer and cli/-adapter modules that sit outside the
%              Tier model — every src/ComplexGitSync/ module is named
%              somewhere in this figure, see docs/DevGuide/architecture.md
%              §2 for the full per-module dependency graph)
%   Figure 2 — GitTree lifecycle state machine
%   Figure 3 — Document hierarchy (ConfigDocument + the ConfigDocumentIOMixin
%              file-I/O split)
%   Figure 3 — Registry model (WorkingGitTree and supporting types)

% ============================================================
%  Shared TikZ style definitions
% ============================================================

\usetikzlibrary{calc}

\tikzset{
  %— UML-style class box with three compartments
  classbox/.style={
    rectangle split,
    rectangle split parts=3,
    rectangle split part fill={blue!10,white,white},
    draw=blue!40,
    thick,
    font=\ttfamily\footnotesize,
    text width=4.8cm,
    align=left,
    inner sep=4pt,
  },
  enumbox/.style={
    rectangle split,
    rectangle split parts=2,
    rectangle split part fill={orange!15,white},
    draw=orange!50,
    thick,
    font=\ttfamily\footnotesize,
    text width=3.8cm,
    align=left,
    inner sep=4pt,
  },
  %— Arrows
  hasarrow/.style={-{Stealth[length=6pt]}, thick},
  inharrow/.style={-{Triangle[open,length=7pt]}, thick},
  comporr/.style={-{Diamond[length=7pt]}, thick},
}

% ============================================================
%  Figure 1 — Class overview by tier
% ============================================================
%
% Same three-tier grouping as figures/three_tier.tex (blue = Tier 3, green =
% Tier 2, orange = Tier 1), but drawn at class granularity: one or two
% headline classbox(es) per tier plus a compact list of every other module
% that tier owns.  The two module groups the book's Tier model never
% described — the cross-cutting Document layer and the Ring-4 cli/ adapter
% package (see docs/DevGuide/architecture.md §1's Tier<->Ring mapping) — get
% their own smaller bands below, so every one of the ~27 modules in
% src/ComplexGitSync/ is named once, either as a classbox or in a module
% list.  Full per-module contracts and the real import graph live in
% docs/DevGuide/architecture.md §2-§3; this figure stays at the zoom level a
% book chapter read start-to-finish needs.

\begin{figure}[p]
\centering
\caption{Class overview by tier.  Headline classes for each of the three
  tiers, plus the cross-cutting Document layer and the \texttt{cli/}
  adapter package that the Tier model never described.  Every module in
  \texttt{src/ComplexGitSync/} is named here, either as a class box or in a
  tier's module list; see \texttt{docs/DevGuide/architecture.md} for the
  full current module-by-module Ring-level dependency graph.}
\label{fig:core-model}

\begin{tikzpicture}[
  font=\sffamily,
  tiertitle/.style={font=\sffamily\bfseries\small},
  modlist/.style={font=\ttfamily\scriptsize, align=left, text width=9.6cm},
  smallnote/.style={font=\ttfamily\scriptsize, align=left, text width=6.4cm},
]

% -------- Tier 3 — Client / API --------
\node[tiertitle] (t3title) at (0,0) {Tier 3 --- Client / API};
\node[classbox, below=0.18cm of t3title.south west, anchor=north west] (client) {
  \textbf{ComplexGitSyncClient}
  \nodepart{two}
  orchestre: Orchestre\\
  git\_runner: GitRunner\\
  state\_store: RuntimeStateStore\\
  registry: WorkingGitTree | None\\
  run\_logger: CommandRunLogger | None
  \nodepart{three}
  configure() / initialise()\\
  pull() / clone() / checkout()\\
  add() / git(tree, cmd)\\
  freeze() / verify()
};
\node[classbox, right=0.5cm of client] (orchestreN) {
  \textbf{Orchestre}
  \nodepart{two}
  git\_tree: GitTree
  \nodepart{three}
  register\_repo(repo)
};
\node[modlist, below=0.3cm of client.south west, anchor=north west] (t3mods) {
  also in \texttt{orchestre.py}: \textit{CommandRunLogger} \textperiodcentered{}
  \textit{RuntimeStateStore} \textperiodcentered{} \textit{LocalGitRegister}
  (\texttt{.lgr} writer) \textperiodcentered{} \textit{SyncLedger}
};

\begin{pgfonlayer}{background}
\node[fill=blue!12, draw=blue!40, thick, rounded corners=6pt, inner sep=0.18cm,
      fit=(t3title)(client)(orchestreN)(t3mods)] (tier3band) {};
\end{pgfonlayer}

% -------- Tier 2 — Actions --------
\node[tiertitle, below=0.5cm of tier3band.south west, anchor=north west] (t2title) {Tier 2 --- Actions};
\node[classbox, below=0.18cm of t2title.south west, anchor=north west] (runner) {
  \textbf{GitRunner}
  \nodepart{two}
  executable: str = ``git''
  \nodepart{three}
  clone() / checkout()\\
  commit() / push() / pull()\\
  create\_branch() / create\_tag()\\
  status\_porcelain() / branch\_tracking\_state()
};
\node[modlist, right=0.6cm of runner] (t2mods) {
  supporting modules (\texttt{Tier 2 --- Actions}):\\
  \textperiodcentered{} \texttt{operations.py} --- leaf/parent-first Git ops\\
  \textperiodcentered{} \texttt{registry.py} --- build/serialize \texttt{WorkingGitTree}\\
  \textperiodcentered{} \texttt{paths.py}, \texttt{discovery.py}\\
  \textperiodcentered{} \texttt{state\_store.py}, \texttt{snapshot\_resolver.py}\\
  \textperiodcentered{} \texttt{ledger\_store.py}
};

\begin{pgfonlayer}{background}
\node[fill=green!10, draw=green!40, thick, rounded corners=6pt, inner sep=0.18cm,
      fit=(t2title)(runner)(t2mods)] (tier2band) {};
\end{pgfonlayer}

\draw[-{Stealth[length=8pt]}, thick, gray]
  (tier3band.south) -- node[right, font=\scriptsize\itshape]{calls (state-gated)} (tier2band.north);

% -------- Tier 1 — Core Data --------
\node[tiertitle, below=0.5cm of tier2band.south west, anchor=north west] (t1title) {Tier 1 --- Core Data};
\node[classbox, below=0.18cm of t1title.south west, anchor=north west] (gittree) {
  \textbf{GitTree} / \textbf{WorkingGitTree}
  \nodepart{two}
  repos: dict[str, GitRepo]\\
  \textit{(WorkingGitTree: dict[str, WorkingRepo],}\\
  \textit{ + lifecycle\_state: TreeLifecycleState)}
  \nodepart{three}
  add\_repo() / force\_repo\_sha()\\
  recompute\_tree\_state()\\
  \textit{WorkingGitTree(GitTree)}
};
\node[classbox, right=0.5cm of gittree] (gitrepo) {
  \textbf{GitRepo} / \textbf{WorkingRepo}
  \nodepart{two}
  project\_name / project\_owner\_name\\
  gitprovider / access\_protocol\\
  commit\_sha
  \nodepart{three}
  resolved\_group\_name\\
  \textit{WorkingRepo(GitRepo)}: repo\_id,\\
  \textit{node\_type, repo\_lifecycle\_state, ...}
};
\node[modlist, below=0.3cm of gittree.south west, anchor=north west] (t1mods) {
  also Tier 1: \texttt{master.py}'s \textit{MasterConfig} (workspace-local
  commit identity) \textperiodcentered{} \textit{RepoAddress} (URL composer,
  \texttt{git\_repo.py})
};

\begin{pgfonlayer}{background}
\node[fill=orange!10, draw=orange!40, thick, rounded corners=6pt, inner sep=0.18cm,
      fit=(t1title)(gittree)(gitrepo)(t1mods)] (tier1band) {};
\end{pgfonlayer}

\draw[-{Stealth[length=8pt]}, thick, gray]
  (tier2band.south) -- node[right, font=\scriptsize\itshape]{reads / mutates} (tier1band.north);

% -------- Cross-cutting: Document layer + cli/ adapter --------
\node[tiertitle, below=0.5cm of tier1band.south west, anchor=north west] (docttitle) {Document layer (cross-cutting)};
\node[smallnote, below=0.18cm of docttitle.south west, anchor=north west] (docmods) {
  \textbf{ConfigDocument} (base) $\to$ \textbf{CgsDocument} / \textbf{GtsDocument},\\
  each also mixing in \textit{ConfigDocumentIOMixin}\\
  (\texttt{config\_document\_io.py}) for \texttt{from/to\_toml|json|yaml}.\\
  Plus \texttt{errors.py}, and the no-Tier-home register/status\\
  internals: \texttt{integrity.py}, \texttt{ledger\_entry.py},\\
  \texttt{status\_render.py}
};

\begin{pgfonlayer}{background}
\node[fill=gray!12, draw=gray!50, thick, rounded corners=6pt, inner sep=0.18cm,
      fit=(docttitle)(docmods)] (docband) {};
\end{pgfonlayer}

\node[tiertitle, right=0.6cm of docband, anchor=north west] (clititle) at (docband.north east) {\texttt{cli/} adapter (Ring 4, outside the Tier model)};
\node[smallnote, below=0.18cm of clititle.south west, anchor=north west] (climods) {
  \texttt{cli/\_\_init\_\_.py} (parser + dispatch)\\
  \texttt{cli/minimalist.py}, \texttt{cli/configuration.py},\\
  \texttt{cli/expert.py} (one module per command group)\\
  \texttt{cli/\_shared.py} (cross-command helpers)\\[2pt]
  \textit{argument/prompt collection only --- calls into}\\
  \textit{Tier 3 above, never touches Git directly}
};

\begin{pgfonlayer}{background}
\node[fill=red!8, draw=red!40, thick, dashed, rounded corners=6pt, inner sep=0.18cm,
      fit=(clititle)(climods)] (cliband) {};
\end{pgfonlayer}

\end{tikzpicture}

\noindent\footnotesize\textit{Not shown: \texttt{\_\_init\_\_.py} (package
re-exports) and \texttt{\_\_main\_\_.py} (entry-point stub) --- package
assembly, not part of any tier or ring.}
\end{figure}

% ============================================================
%  Figure 2 — GitTree lifecycle states
% ============================================================

\begin{figure}[ht]
\centering
\caption{GitTree lifecycle state machine.  The simplified user-facing lifecycle
  is \texttt{initialise(.cgs)} $\rightarrow$ clone $\rightarrow$ \texttt{READY};
  \texttt{initialise(.gts)} $\rightarrow$ restore $\rightarrow$ \texttt{READY};
  \texttt{pull(.cgs|.gts)} $\rightarrow$ resync $\rightarrow$ \texttt{READY}.
  Synchronized Git operations (\texttt{checkout},
  \texttt{add}, \texttt{git("commit")}, \texttt{git("push")},
  \texttt{git("tag")},
  \texttt{freeze}) are gated on \texttt{READY}.}
\label{fig:gittree-lifecycle}

\begin{tikzpicture}[
  node distance=1.4cm and 2.0cm,
  state/.style={rectangle, rounded corners=4pt, draw=blue!50, fill=blue!8,
                font=\ttfamily\small, minimum width=2.6cm, minimum height=0.8cm,
                inner sep=4pt},
  sidestate/.style={rectangle, rounded corners=4pt, draw=orange!60, fill=orange!10,
                    font=\ttfamily\small, minimum width=2.6cm, minimum height=0.8cm,
                    inner sep=4pt},
  trans/.style={-{Stealth[length=6pt]}, thick},
]

\node[state] (unloaded) {UNLOADED};
\node[state, right=of unloaded] (declared)  {DECLARED};
\node[state, right=of declared] (pending)   {PENDING};
\node[state, right=of pending]  (ready)     {READY};

\node[sidestate, below=1.2cm of declared] (partial) {PARTIAL};
\node[sidestate, below=1.2cm of pending]  (error)   {ERROR};

\draw[trans] (unloaded) -- node[above, font=\scriptsize]{load(.cgs)} (declared);
\draw[trans] (declared) -- node[above, font=\scriptsize]{discover nested} (pending);
\draw[trans] (pending)  -- node[above, font=\scriptsize]{clone/validate} (ready);
\draw[trans, dashed] (declared.south) -- (partial.north);
\draw[trans, dashed] (pending.south)  -- (error.north);
\draw[trans, bend right=20] (ready.north) to node[above, font=\scriptsize]{freeze / next .gts id + .lgr update} (declared.north);

\end{tikzpicture}
\end{figure}

% ============================================================
%  Figure 3 — Document class hierarchy
% ============================================================

\begin{figure}[ht]
\centering
\caption{Document class hierarchy.  \texttt{ConfigDocument} (\texttt{config\_document.py})
  is the abstract, I/O-free base; \texttt{ConfigDocumentIOMixin}
  (\texttt{config\_document\_io.py}) supplies the file-based
  \texttt{from/to\_toml|json|yaml} methods.  Each concrete document type
  combines both via ordinary multiple inheritance.}
\label{fig:document-hierarchy}

\begin{tikzpicture}[node distance=1.2cm and 1.8cm]

\node[classbox] (base) {
  \textbf{ConfigDocument}
  \nodepart{two}
  config\_document.py\\
  \_data: dict
  \nodepart{three}
  read(key) / get(key)\\
  validate()\\
  to\_dict() / from\_dict(data)
};

\node[classbox, right=2.4cm of base] (iomixin) {
  \textbf{ConfigDocumentIOMixin}
  \nodepart{two}
  config\_document\_io.py
  \nodepart{three}
  from\_toml / from\_json / from\_yaml\\
  to\_toml / to\_json / to\_yaml\\
  print()
};

\node[classbox, below left=1.6cm and 0.6cm of base] (cgs) {
  \textbf{CgsDocument}
  \nodepart{two}
  cgs\_format.py\\
  DOCUMENT\_KIND = ``cgs''
  \nodepart{three}
  project\_name\\
  default\_branch\\
  repos\\
  normalize authoring syntax\\
  to\_git\_tree() / from\_git\_tree()\\
  to\_authoring\_dict() / to\_toml()\\
  runtime\_setting(key)
};

\node[classbox, below=1.6cm of base] (gts) {
  \textbf{GtsDocument}
  \nodepart{two}
  gts\_document.py\\
  DOCUMENT\_KIND = ``gts''
  \nodepart{three}
  schema\_version\\
  snapshot\_hash\\
  compute\_snapshot\_hash()\\
  ensure\_snapshot\_hash()\\
  lifecycle\_state\\
  is\_ready\\
  repo\_states
};

\draw[inharrow] (cgs.north) -- (base.south west);
\draw[inharrow] (gts.north) -- (base.south);
\draw[inharrow, dashed] (cgs.north east) -- ($(iomixin.south)+(-0.6,0)$);
\draw[inharrow, dashed] (gts.east) -- (iomixin.south);

\end{tikzpicture}
\end{figure}

% ============================================================
%  Figure 4 — Registry model
% ============================================================

\begin{figure}[ht]
\centering
\caption{Registry model.  \texttt{WorkingGitTree} is the authoritative
  in-memory graph; \texttt{WorkingRepo} \emph{extends} \texttt{GitRepo}
  (Figure~\ref{fig:core-model}) rather than duplicating its identity
  fields, adding only the mutable tree-position and synchronisation state
  below.}
\label{fig:registry-model}

\begin{tikzpicture}[node distance=1.2cm and 2.0cm]

% WorkingGitTree
\node[classbox] (registry) {
  \textbf{WorkingGitTree}
  \nodepart{two}
  repos: dict[str, WorkingRepo]\\
  lifecycle\_state: TreeLifecycleState
  \nodepart{three}
  add(entry) / get(id) / values()\\
  children\_of(parent\_id)\\
  is\_complete() / is\_ready()\\
  recompute\_tree\_state()
};

% WorkingRepo
\node[classbox, right=2.2cm of registry] (entry) {
  \textbf{WorkingRepo} \textit{(extends GitRepo)}
  \nodepart{two}
  repo\_id / name / node\_type\\
  parent\_id / absolute\_path\\
  current\_ref\_kind / \_name\\
  target\_ref\_kind / \_name\\
  resolved\_ref\_kind / \_name\\
  repo\_lifecycle\_state\\
  sync\_state / discovery\_state\\
  \textit{(+ inherited: project\_name,}\\
  \textit{gitprovider, commit\_sha, ...)}
  \nodepart{three}
  (mutable dataclass)
};

% ProjectTreeState
\node[classbox, below=1.6cm of registry] (treestate) {
  \textbf{ProjectTreeState}
  \nodepart{two}
  lifecycle\_state: TreeLifecycleState\\
  is\_ready: bool\\
  registry\_complete: bool
  \nodepart{three}
  (frozen dataclass)
};

% Enum: NodeType
\node[enumbox, below left=1.6cm and 0.2cm of entry] (nodetype) {
  \textbf{NodeType} \textit{(enum)}
  \nodepart{two}
  ROOT / PARENT / LEAF
};

% Enum: TreeLifecycleState
\node[enumbox, right=1.8cm of treestate] (tlc) {
  \textbf{TreeLifecycleState}
  \nodepart{two}
  UNLOADED / DECLARED\\
  PENDING / READY\\
  PARTIAL / ERROR
};

% Enum: RepoLifecycleState
\node[enumbox, below=1.2cm of nodetype] (rlc) {
  \textbf{RepoLifecycleState}
  \nodepart{two}
  DECLARED / PENDING\\
  READY / FALLBACK\_READY\\
  MISSING / ERROR
};

% Arrows
\draw[comporr] (registry.east) -- node[above, font=\scriptsize]{0..*} (entry.west);
\draw[hasarrow] (registry.south) -- (treestate.north);
\draw[inharrow, dashed] (entry.north) -- ++(0,1.0) node[above, font=\scriptsize\itshape]{GitRepo (Fig.~\ref{fig:core-model})};
\draw[hasarrow] (entry.south west) -- (nodetype.north east);
\draw[hasarrow] (entry.south) -- (rlc.north);
\draw[hasarrow] (registry.south east) -- (tlc.north west);

\end{tikzpicture}
\end{figure}
